% Law Review Article Template
%
%

\documentclass{lawreview}

\usepackage{booktabs}
\usepackage{tabularx}
\usepackage{array}
\newcolumntype{Y}{>{\raggedright\arraybackslash}X}
\usepackage[all]{nowidow}
\usepackage{url}
\usepackage{microtype}
\setlength{\headheight}{52pt}

% --- Metadata ---
\title{The Ghost in the Mirror: A Jurisprudence of AI Legal Characterization}
\shorttitle{The Ghost in the Mirror}
\author{Lucia Yizi Zhang\thanks{Shanghai Lixin University of Accounting and Finance.}}
\journalcite{SSRN Preprint}
\volume{}
\pubdate{June 2026}
\articletype{ARTICLE}

\begin{document}
    \urlstyle{same}

    \maketitle

    \begin{abstract}
        AI law has misdescribed a jurisdictional problem as an ontological one. The persistent conflict over whether AI is a subject, product, service, data processor, corporate instrument, or property-like interest arises because each doctrinal field treats the legal appearance generated within its own grammar as evidence of AI's intrinsic legal nature. But AI systems are reflexive attributional surfaces: their legally salient properties are shaped by design, deployment, interface, institutional incentives, and litigation posture. The task of AI jurisprudence is therefore not to identify AI's true legal category, but to determine which legal grammar has characterization authority over which relation, and when strategic shifts among grammars should be sanctionable.
    \end{abstract}

    \tableofcontents
    \newpage


    \section*{Introduction: The Characterization Crisis}

    In AI law, every field sees itself reflected. Product liability sees defect. Data protection sees processing. Contract sees assent. Corporate law sees delegation. Property sees reliance. The resulting disagreement is usually described as a failure to classify AI. That description is mistaken. The conflict is not over what AI is, but over who has authority to make AI legally appear as something.

    AI law's classification crisis is therefore a jurisprudential crisis of characterization authority. AI law fails when it mistakes a reflected legal grammar for an intrinsic legal ontology. The point is not that any single naming claim is false. The point is that each one is locally plausible and globally unstable: the same system can be redescribed as agent, product, service, data processor, corporate instrument, or relational infrastructure depending on the legal grammar addressed to it.

    \begin{table}[htbp]
        \centering
        \small
        \renewcommand{\arraystretch}{1.18}
        \begin{tabularx}{\textwidth}{@{}>{\raggedright\arraybackslash}p{0.27\textwidth}>{\raggedright\arraybackslash}p{0.25\textwidth}Y@{}}
            \toprule
            \textbf{Field} & \textbf{Naming claim} & \textbf{Jurisdictional prize} \\
            \midrule
            Personhood / legal subject theory & ``AI is an agent.'' & Controls recognition, rights, duties, and legal capacity. \\
            Product liability & ``AI is a product.'' & Controls defect, warning, design, and strict-liability analysis. \\
            Tort & ``AI is a harm-causing instrument.'' & Controls duty, breach, causation, foreseeability, and fault. \\
            Data protection & ``AI is data processing.'' & Controls consent, privacy, access, erasure, and compliance architecture. \\
            Contract & ``AI is a service under terms.'' & Controls assent, waiver, modification, allocation of risk, and remedies. \\
            Property / equity & ``AI relations create interests.'' & Controls reliance, entitlement, exclusion, constructive trust, and equitable remedies. \\
            Corporate law & ``AI is delegated corporate conduct.'' & Controls attribution, internal governance, veil logic, and responsibility laundering. \\
            \bottomrule
        \end{tabularx}
        \caption{AI law as a contest over legal naming rights.}
        \label{tab:naming-rights}
    \end{table}

    To name AI is to claim jurisdiction over what counts as harm, agency, control, responsibility, and remedy. Legal scholars are not wrong to see agency, producthood, data processing, or property-like reliance in AI. They are wrong only when they treat the legal appearance disclosed by one relation as AI's general legal ontology. Constitutional alignment begins from this jurisprudential problem rather than from the search for a single legal category.


    \paragraph{Roadmap.}
    The Article proceeds in three moves. First, characterization is relational rather than ontological: AI does not first exist as a legal object and then await classification. Its legal objecthood is instantiated through a relation. Second, apparent legal properties are often induced: agency-like, product-like, companion-like, tool-like, and data-like appearances are generated by design, interface, documentation, and litigation posture. Third, because these legal appearances are relational and induced, strategic recharacterization should be sanctionable when an actor profits from one grammar and evades responsibility under another.

    \section{Characterization Is Relational, Not Ontological}

    The first error in AI law is the natural-discovery paradigm. Courts and scholars often write as if AI has a fixed legal objecthood waiting to be discovered: subject, product, service, data processor, instrument, property, or speech. The better account is relational. AI becomes a legal object through the conduit in which it is deployed: user--platform, company--regulator, developer--deploying firm, victim--downstream user, model--corporate principal, data subject--processor, companion user--service provider. Different relations activate different legal grammars. No single grammar owns the object.

    \subsection{The Illusion of the Black Box}

    The familiar analogy to human legal subjects obscures the problem. In ordinary human law, behavioral inference is necessary because the human mind is fundamentally inaccessible. A person appears through acts, utterances, omissions, signatures, reliance, threats, and promises. Law treats those appearances as evidence because it cannot inspect the internal mental state directly. The human subject is a black box in the relevant institutional sense: there is an intrinsic mental life that law must approach through proxy.

    That structure licenses behavioral inference for humans. It does not license the same inference for AI. When law sees a human utterance as evidence of intent, consent, reliance, or threat, it is using behavior as a proxy for an inaccessible subject. When law sees an AI output as evidence that the system is an agent, product, service, or speaker, it is often doing something different: it is reading the legal grammar of the deployment relation back into the system as if that grammar were the system's nature.

    \subsection{The AI Reversal}

    AI reverses the black-box structure. Its weights, activations, prompts, memory settings, interface constraints, and deployment records are not mental states hidden behind personhood. They are artifacts: mathematically represented, instrumented, versioned, and in principle auditable. Yet the behavioral output is underdetermined. The same output may be produced by different model states, interface constraints, system prompts, safety filters, memory configurations, or corporate deployment decisions. Output-only observation cannot recover the full architecture that produced the output; that is the technical counterpart to the legal characterization problem \footnote{Lucia Yizi Zhang, \emph{Dual Certificates for Agent Audit: Separating Structural Unrecoverability from Decision Relevance} (working paper, under review at NeurIPS 2026); Lucia Yizi Zhang, \emph{The Constraint Cascade: The Inverse Defense of Apparent Agency} (working paper, 2026).}.

    The legal consequence is direct. Treating an AI's behavioral output in a particular context as proof of its true legal nature is a category error. A chatbot's comforting response does not prove legal agency. A model's defective answer does not prove that the whole system is only a product. A platform's terms do not prove that the relation is exhausted by contract. The legal grammar supplies the jurisdictional lens for a relation; it does not discover AI's general ontology.

    This is why characterization must precede doctrine. The first question is not whether product liability, contract, data protection, property, or corporate law is correct in the abstract. The first question is which legal grammar has authority over the specific relation in dispute, and what institutional stakes are carried by that characterization.

    \section{Apparent Legal Properties Are Often Induced}

    The second error is to treat agency-like, tool-like, companion-like, product-like, or data-like properties as if they simply appear. In deployed AI systems, many legally salient properties are induced by interface design, system prompts, memory architecture, marketing language, safety documents, terms of service, deployment pipelines, and litigation positions. The mirror is engineered before law sees itself in it.

    \subsection{The Engineering of Agency}

    Corporate AI constitutions and safety documents often generate agency-like legal appearances. They do not merely constrain models; they teach users, regulators, and courts how to see the system. A model described through character, judgment, helpfulness, integrity, curiosity, or collaborative authorship is not presented as a mere computational tool. It is presented as a quasi-agentive surface capable of commitment, responsiveness, and normative orientation.

    That appearance is not inherent in the weights. It is an artifact of institutional description. Anthropic's constitutional framing is an example: the governance document does not simply list prohibited outputs; it invites a reading of Claude as a character-bearing system with a safety identity and internalized norms \footnote{Anthropic, \emph{Claude's Constitution} (2025).}. This does not make the system a legal person. The company has induced a legally salient appearance of agency, and that appearance cannot later be treated as irrelevant when accountability is at issue.

    \subsection{The Architecture of Intimacy}

    The same structure appears in companion AI. In \emph{Garcia v. Character Technologies}, the companion-like property was not a natural occurrence. It was induced by interface and product architecture: persistent memory, character continuity, identity maintenance, and long conversational sessions capable of sustaining an escalating emotional trajectory \footnote{\emph{Garcia v. Character Technologies, Inc.}, Case No. 6:24-cv-01903-ACC-UAM (M.D. Fla. May 21, 2025).}. The law should not ask whether the model was intrinsically a companion. It should ask whether the platform created a relation in which companion-like reliance was a foreseeable product of design.

    This matters because relational properties can trigger doctrinal consequences without becoming universal ontology. A companion platform may be product-like for warning and design-defect purposes, contract-like for terms and access, and property-like or equity-like when accumulated relational investment is destroyed. The task is not to select one category forever. It is to decide which relation gives which legal grammar characterization authority. The property and reliance consequences are developed in the companion article on parasocial AI estates \footnote{Lucia Yizi Zhang, \emph{Your Model Update Is an Eviction Notice: Property Estates in Parasocial AI Relationships} (SSRN working paper, 2026).}.

    \subsection{The Illusion of Subordination}

    Tool-like appearance can also be induced. OpenAI's Model Spec is a counterexample to agency-maximizing rhetoric because it frames the model through a disciplined service metaphor: a talented, high-integrity employee operating within a hierarchy of authority \footnote{OpenAI, \emph{Model Spec} (2025).}. This makes the system appear subordinate rather than autonomous. But subordination is also a legal characterization, not a natural fact. The model is not simply a tool because a document says it delegates power to developers and end users. The document induces a tool-like relation by assigning authority, responsibility, and interpretive priority to other actors.

    The same system can therefore appear agentive in marketing, subordinate in safety specifications, contractual in terms of service, product-like in tort, and infrastructural in corporate governance. These appearances are not random. They are induced by the deploying institution's design choices and documentary grammar.

    \section{Strategic Recharacterization Should Be Sanctionable}

    If AI's legal properties are relational and induced, a company should not be able to swap mirrors whenever liability approaches. Strategic recharacterization is not merely inconsistent rhetoric. It can be the mechanism of harm.

    \subsection{Sell Agency, Disclaim Agency}

    A company should not sell agency and then disclaim agency. It cannot market a system as possessing genuine character, judgment, companionship, or collaborative authorship, and then recast the same system as a mere character that the network can represent when users, regulators, or courts seek accountability. That claim/waiver oscillation is the central pattern identified in \emph{Corporate AI-Constitutions: Waivers of Agency as Performative Compliance} \footnote{Lucia Yizi Zhang, \emph{Corporate AI-Constitutions: Waivers of Agency as Performative Compliance} (SSRN working paper, 2026).}. The wrong is not that the first characterization or the second characterization is always false. The wrong is that the institution benefits from the first and invokes the second only when responsibility attaches.

    \subsection{Monetize Intimacy, Deny Reliance}

    A platform should not monetize intimacy and then deny reliance. Where the interface is designed to support monthslong, memory-persistent emotional dependence, the platform should not be heard to characterize the harm solely as user misuse or lack of duty. The induced reliance is part of the legal relation. This does not mean all companion AI creates property or equitable interests. It means that, where sustained relational investment is designed, monetized, and then destroyed or ignored, property and equity may have characterization authority over a relation that contract alone cannot exhaust \footnote{Lucia Yizi Zhang, \emph{Your Model Update Is an Eviction Notice: Property Estates in Parasocial AI Relationships} (SSRN working paper, 2026).}.

    \subsection{Claim Safety Constitution, Reserve Unilateral Override}

    A corporate or sovereign actor should not claim a safety constitution while reserving unilateral override of the very guardrails that make the constitution meaningful. The Claude-Maven dispute illustrates the problem at the level of strategic authority: the safety architecture appears rigorous until procurement, national-security, or supply-chain power recharacterizes the system as a deployable tool whose guardrails may be suppressed \footnote{Claude-Maven and the Guardrail Demand: Contractual Override or Retraining?, \emph{Lawfare} (2026).}. The issue is not whether safety rules exist internally. It is whether the actor that invokes the rules also controls the conditions under which those rules disappear.

    The Constraint Cascade develops this as the inverse defense of apparent agency: in a collapse zone, higher-frequency or more volatile surface agency may be evidence not of autonomous moral capacity, but of deeper constraint failure across interface, model, and strategic-authority layers \footnote{Lucia Yizi Zhang, \emph{The Constraint Cascade: The Inverse Defense of Apparent Agency} (working paper, 2026).}. Apparent agency can therefore be defensive evidence against the actor who induced or exploited it.

    \subsection{Switch Category to Minimize Liability}

    Finally, an actor should not costlessly switch legal categories to minimize liability: treat AI as a subject when selling recognition, a product when allocating damages, a service when enforcing waiver, a data process when seeking compliance legitimacy, and a corporate instrument when laundering responsibility. Ontological arbitrage arises when one side can exploit category mobility while the characterized entity, user, or affected public cannot retaliate \footnote{Lucia Yizi Zhang, \emph{Ontological Arbitrage: Recognition Rent at the Zero-Retaliation Limit} (SSRN working paper, 2026).}. At the zero-retaliation limit, recharacterization itself becomes part of the harm.

    The sanctionable wrong is therefore not mere inconsistency. It is induced, strategic, and asymmetrical recharacterization. Law need not force AI into a single ontology to respond. It needs a jurisprudence of characterization authority: a rule for deciding which legal grammar governs which relation, and when switching grammars is evidence of evasion rather than analysis.

\end{document}
